
\medskip
\phantomsection\label{lem:sga3-VIB-5.6.1.0}%
\noindent\textbf{Lemma 5.6.1.0.---}%
\footnote{\label{note:sga3-VIB-39}Editor's note: we have made
explicit this lemma, which is used in the proof of \hyperref[lem:sga3-VIB-5.6.1]{Lemma 5.6.1}.}
\emph{Let \(f:X\to S\) be a morphism locally of finite
presentation, let \(s\in S\), and let \(x\in X_s\). Set
\(n=\dim_x(X_s)\), and let
\(q:\mathbb A^n_S\to S\) be the structure morphism.}

\emph{Suppose that a quasi-finite morphism
\(u:X\to\mathbb A^n_S\) is given such that \(f=q\circ u\), and
suppose that \(f\) is universally open at the generic point \(z\)
of an irreducible component \(\mathfrak z\) of \(X_s\), containing
\(x\) and of dimension \(n\). Then \(u\) is universally open at
the point \(x\in X\).}

First observe that
\[
  (\dagger)\qquad
  \text{\emph{the hypotheses are preserved under every base change
  \(\pi:S'\to S\) covering \(s\)}}
\]
(that is, such that \(\pi^{-1}(s)\ne\varnothing\)).
Indeed, let \(S'\to S\) be such a morphism, let
\begin{center}
  \phantomsection\label{diag:sga3-VIB-5.6.1.0-base-change}%
  \begin{tikzcd}[row sep=large,column sep=large]
    X'
      \arrow[r,"u'"]
      \arrow[dr,"f'"'] &
    \mathbb A^n_{S'} \arrow[d,"q'"] \\
    {} &
    S'
  \end{tikzcd}
\end{center}
be the diagram obtained by base change, and let \(s'\) be a point
of \(S'\) above \(s\) and \(x'\) a point of \(X'_{s'}\) above
\(x\). Since
\[
  X'_{s'}=X_s\otimes_{\kappa(s)}\kappa(s'),
\]
the lifting of generizations and the invariance of dimension under
field extension (EGA~IV\(_2\), 2.3.4\textup{(i)} and 4.1.4) show
that \(x'\) is contained in an irreducible component
\(\mathfrak z'\) of \(X'_{s'}\) whose generic point \(z'\) lies
above \(z\), and
\[
  n\leq\dim\mathfrak z'
    \leq\dim_{x'}X'_{s'}\leq n.
\]
Thus
\[
  \dim\mathfrak z'=n=\dim_{x'}X'_{s'}.
\]
Since \(f\) is universally open at \(z\), \(f'\) is universally
open at \(z'\).%
\footnote[1]{\label{note:sga3-VIB-5.6.1.0-generic-points}Translator's
note: the French source prints \(x\) and
\(x'\), respectively, in this sentence. The stated hypothesis and
the immediately preceding choice of generic points require \(z\)
and \(z'\).}

Set \(Y=\mathbb A^n_S\). By EGA~IV\(_3\), 14.3.3.1\textup{(i)},
to prove that \(u\) is universally open at \(x\), it is enough to
prove that, for every integer \(r\geq0\) and every point \(x'\) of
\[
  X'=X[T_1,\ldots,T_r]
\]
above \(x\), the morphism
\[
  u':X'\longrightarrow
  Y'=\mathbb A^n_S[T_1,\ldots,T_r]
\]
is open at \(x'\). Now set
\[
  S'=S[T_1,\ldots,T_r],
\]
and let \(q'\) be the projection
\[
  Y'\simeq\mathbb A^n_{S'}\longrightarrow S'.
\]
This is the situation obtained by the base change \(S'\to S\).
Thus, by the preceding observation, we are reduced to showing that
\(u\) is open at \(x\).

Set \(y=u(x)\). Since \(u\) is locally of finite presentation
(because \(f\) and \(q\) are; cf. EGA~IV\(_1\), 1.4.3),
EGA~IV\(_1\), 1.10.3 shows that it is enough to prove
\[
  u\bigl(\operatorname{Spec}\mathcal O_{X,x}\bigr)
  =\operatorname{Spec}\mathcal O_{Y,y}.
\]
For this purpose we may assume that \(S\) is affine and integral.
Let \(\pi:S'\to S\) be its normalization, and denote by
\(u':X'\to Y'\) the morphism obtained from \(u\) by base change.
Since
\[
  \pi_Y:Y'\longrightarrow Y
\]
is integral and surjective, one has
\[
  \operatorname{Spec}\mathcal O_{Y,y}
  =
  \bigcup_{y'}
  \pi_Y\bigl(\operatorname{Spec}\mathcal O_{Y',y'}\bigr),
\]
where the union ranges over all points \(y'\) of \(Y'\) above
\(y\). It is therefore enough to show that, for every such
\(y'\) and every \(x'\in X'_{y'}\) above \(x\),
\[
  u'\bigl(\operatorname{Spec}\mathcal O_{X',x'}\bigr)
  =\operatorname{Spec}\mathcal O_{Y',y'}.
\]
Since the hypotheses are preserved under the base change
\(\pi:S'\to S\), we are thereby reduced to the case of \(S'\);
that is, we may assume that \(S\) is integral and normal.

The hypotheses on \(f\) now imply, by EGA~IV\(_3\), 14.3.13, that
there is an irreducible component \(Z\) of \(X\), containing
\(\mathfrak z\) (and hence \(x\)), dominating \(S\), and such that
\[
  \dim_z(Z_s)=n=\dim(Z_\eta),
\]
where \(\eta\) is the generic point of \(S\). Let \(\xi\) be the
generic point of \(Z\). Since \(u\), and hence also
\[
  u_\eta:Z_\eta\longrightarrow\mathbb A^n_\eta,
\]
is quasi-finite, the closure in \(\mathbb A^n_\eta\) of the point
\[
  u_\eta(\xi)=u(\xi)
\]
has dimension \(n\). Thus \(u(\xi)\) is the generic point of
\(\mathbb A^n_\eta\), which is also the generic point of
\(Y=\mathbb A^n_S\). Consequently, if \(g\) denotes the
restriction of \(u\) to \(Z\), then
\[
  g:Z\longrightarrow Y
\]
is quasi-finite and dominant. Since \(Y\) is normal, \hyperref[lem:sga3-VIB-5.6.0]{Lemma 5.6.0}
shows that \(g\) is open. Hence \(u\) is open at every point of
\(Z\), and in particular at \(x\). Consequently,
\[
  u\bigl(\operatorname{Spec}\mathcal O_{X,x}\bigr)
  =\operatorname{Spec}\mathcal O_{Y,y},
\]
which completes the proof of \hyperref[lem:sga3-VIB-5.6.1.0]{Lemma 5.6.1.0}.

The following lemma advantageously replaces
EGA~IV\(_3\), 14.5.10,%
\footnote{\label{note:sga3-VIB-40}Editor's note: we have corrected
the original, which indicated 19.5.10.}
in that it is independent of noetherian hypotheses.
